\phantomsection
\subsection*{M2B. Pyment}
\addcontentsline{toc}{subsection}{M2B. Pyment}

\textit{Pyment é um melomel elaborado com uvas de vinho ou de mesa (geralmente a partir de suco ou mosto). Assim como ocorre com o vinho, o pyment pode ser tinto, branco ou rosé.}

\textbf{Impressão Geral}: Assemelha-se a um M1 Traditional Mead, porém com um caráter perceptível de vinho. Uma ampla gama de resultados é possível, dependendo das escolhas do hidromel base e dos ingredientes adicionais, mas o produto final deve ser equilibrado, agradável e prazeroso. O caráter da fruta não deve dominar completamente o hidromel.

\textbf{Aroma}: Os aromas da uva podem ser percebidos como de um suco doce e fresco até um vinho reserva envelhecido, dependendo da força alcoólica e do dulçor. Se variedades de uva forem declaradas, o caráter varietal deve estar perceptível. Pyments costumam apresentar uma sugestão de acidez ou de taninos no aroma.

O caráter do mel pode variar de sutil a intenso, dependendo da força e dulçor, podendo expressar o néctar floral típico de eventuais variedades de mel declaradas. O buquê de fermentação deve ser equilibrado. Hidroméis mais fortes podem apresentar uma leve nota alcoólica picante, e os carbonatados tendem a expressar mais aroma.

Os aromas provenientes da uva devem complementar e realçar o hidromel base, não sobrepujá-lo.

\textbf{Aparência}: Limpidez de boa a brilhante; quanto maior a limpidez, mais desejável. A presença de bolhas, espuma, colarinho ou efervescência depende do nível de carbonatação declarado. A cor pode variar de palha claro a vermelho-púrpura profundo, dependendo da variedade de uva e de mel utilizadas. A cor deve ser característica da variedade ou tipo de uva usada, embora exemplares com variedades de uva branca também possam adquirir cor derivada da variedade de mel. Quanto maior a vivacidade e o brilho da cor, mais desejável.

\textbf{Sabor}: Semelhante ao Aroma quanto a característica e intensidade de fruta e mel, se tornando mais proeminente em versões mais alcoólicas e doces. O caráter da fruta deve ser perceptível e equilibrado com o nível de dulçor do hidromel. Algumas frutas podem conferir gostos básicos (amargo, doce, ácido) além de suas características típicas. Muitas uvas acrescentam acidez e taninos naturais, que precisam estar em equilíbrio com o hidromel base.

Dulçor residual e final de acordo com o nível de dulçor declarado. A acidez deve ser equilibrada, macia ou vibrante, mas não agressiva. O caráter de uva e de fermentação típico de vinhos também é aceitável em Pyment. Taninos podem fazer um hidromel suave parecer mais seco. Características ásperas, desagradáveis ou excessivamente sulfurosas ou com notas de levedura/autólise são indesejáveis.

\textbf{Sensação na Boca}: A carbonatação é indicada pelo nível declarado, de tranquilo a espumante. O corpo geralmente aumenta com o dulçor e o teor alcoólico, tipicamente de médio-leve a cheio. A percepção de álcool acompanha o teor alcoólico declarado, variando de ausente a perceptível e com sensação de aquecimento. Os taninos provenientes das uvas ou das cascas (particularmente em variedades tintas) conferem algum corpo e uma sensação adstringente, que não deve ser excessiva.

\textbf{Ingredientes}: Os mesmos ingredientes de um M1 Traditional Mead, com a adição de pelo menos um tipo de uva ou obtido a partir da mistura de hidromel e vinho.

\textbf{Comentários}: Um Pyment que também contenha outra fruta deve ser inscrito como M2E Other Fruit Mead. Pyments com especiarias (Hippocras) devem ser inscritos como M3C Fruit and Spice Mead. Hidroméis com frutas e outros ingredientes que não forem frutas provavelmente devem ser inscritos na categoria M4F (Experimental Mead).

\textbf{Instruções para Inscrição}: Os participantes devem especificar os níveis de dulçor, carbonatação e teor alcoólico. Variedades de mel podem ser declaradas. A fruta deve obrigatoriamente ser declarada. Se variedades de uva forem especificadas, espera-se que o caráter varietal correspondente esteja presente.

\textbf{Exemplos Comerciais}: Aurum Verdelho Pyment, Hammen Family Orchard Pyment, Meadow Vista Ostara, Mount Stirling Red Pyment Mead, Remeadies Sour Grapes.
